\documentclass[12pt,a4paper]{article}

% Packages
\usepackage[utf8]{inputenc}
\usepackage[T1]{fontenc}
\usepackage{amsmath,amssymb,amsfonts}
\usepackage{graphicx}
\usepackage{booktabs}
\usepackage{longtable}
\usepackage{multirow}
\usepackage{array}
\usepackage{float}
\usepackage{hyperref}
\usepackage[margin=1in]{geometry}
\usepackage{natbib}
\usepackage{setspace}
\usepackage{caption}
\usepackage{subcaption}
\usepackage{xcolor}
\usepackage{enumitem}

% Custom commands
\newcommand{\bps}{\text{ bps}}
\newcommand{\ETH}{\text{ETH}}
\newcommand{\USDC}{\text{USDC}}

% Title
\title{\textbf{Execution Quality in Order Flow Auctions: \\
Evidence from CoW Protocol and 1inch Fusion}\\[1em]
\large A Comprehensive Empirical Analysis of Solver Behavior, \\
Price Momentum, and Market Microstructure}

\author{Research Report}
\date{January 2026}

\begin{document}

\maketitle

\begin{abstract}
We conduct a comprehensive empirical analysis of execution quality in Order Flow Auctions (OFAs), comparing CoW Protocol and 1inch Fusion across 69,661 trades over 32 days in Summer 2025. Using forward and backward markouts relative to Binance prices, we document several novel findings. First, we discover that \textbf{CoW Protocol's batch auction mechanism provides measurable execution benefits} (+0.38--0.50 bps) relative to 1inch Fusion, even for trades using the same settlement infrastructure. Second, we identify \textbf{distinct CoW solvers outperform}: trades routed through identifiable solver addresses achieve +3.57 bps versus +3.07 bps for 1inch settlement trades (t=4.42, p$<$0.001), demonstrating the \textbf{value of solver competition}. Third, we show that price momentum---not adverse selection---is the primary driver of markout distribution, with backward-forward correlations of 0.54--0.62 for major token pairs. Fourth, we find that higher volatility correlates with \textit{better} markouts, contradicting Loss-Versus-Rebalancing (LVR) theory, but explain this through sample composition: 79\% of trades are sells in predominantly rising markets. Fifth, we demonstrate that solvers cannot profitably time CEX-DEX latency arbitrage---``favorable'' timing produces worse markouts due to momentum persistence. Sixth, we document \textbf{striking directional asymmetry across token pairs}: ETH$\rightarrow$Stablecoin trades achieve +0.24\% markout with 77\% win rate, while Stablecoin$\rightarrow$ETH trades show -0.40\% markout with only 27\% win rate---a 64 basis point spread that dwarfs platform effects. Notably, \textbf{competition matters most in adverse conditions}: during falling markets, the CoW advantage increases to +1.18 bps. Our extended analyses reveal that users capture only 35\% of available price surplus, execution consistency is low (CV $>$ 2.5), approximately 14\% of order flow shows potentially toxic patterns, and market regime is the dominant factor explaining 70\% of markout variance---with a 22 bps spread between best and worst regimes. Our findings reveal that while mechanism design and solver competition contribute to execution quality, \textbf{trade direction relative to market trend is the dominant determinant of user outcomes}.

\vspace{1em}
\noindent\textbf{Keywords:} Order Flow Auctions, Decentralized Exchanges, Market Microstructure, MEV, Execution Quality, Solvers, CoW Protocol, 1inch Fusion, Batch Auctions, Mechanism Design

\vspace{1em}
\noindent\textbf{JEL Classification:} G14, G23, D47, L12
\end{abstract}

\newpage
\tableofcontents
\newpage

%==============================================================================
\section{Introduction}
%==============================================================================

The emergence of Order Flow Auctions (OFAs) represents a fundamental shift in how retail traders access liquidity in decentralized finance (DeFi). Unlike traditional Automated Market Makers (AMMs) where users execute directly against liquidity pools, OFAs introduce competitive intermediaries---called \textit{solvers}---who bid for the right to fill user orders. This mechanism promises to improve execution quality through competition, route optimization, and protection from Maximal Extractable Value (MEV) extraction.

Despite their growing importance---with platforms like CoW Protocol and 1inch Fusion collectively processing billions in volume---empirical evidence on OFA execution quality remains limited. Existing research has focused primarily on theoretical mechanism design \citep{flashbots2023ofa1, flashbots2023ofa2} or narrow empirical questions \citep{bachu2024, yuminaga2025}. Critical questions remain unanswered: Do solvers strategically time order execution based on CEX price movements? How does volatility affect execution quality in OFAs compared to theoretical predictions? Does the batch auction mechanism provide value beyond what competitive solvers alone deliver?

This paper provides the first comprehensive empirical analysis addressing these questions. We analyze 69,661 trades across CoW Protocol and 1inch Fusion over four sampling periods spanning June--September 2025. Our methodology employs forward and backward markouts---measuring price changes from -60 seconds to +60 seconds around execution---using Binance mid-prices as the benchmark.

Our analysis yields five main findings:

\begin{enumerate}[leftmargin=*]
    \item \textbf{CoW Protocol's batch auction mechanism provides execution benefits.} We find that CoW Protocol outperforms 1inch Fusion by +0.38--0.50 bps on average. Importantly, our data reveals that 1inch Fusion trades route through the GPv2Settlement contract (address 0x9008D19f58...), preventing resolver-level analysis on that platform. However, CoW Protocol's six identifiable solver addresses enable direct measurement of solver-level competition effects.

    \item \textbf{Distinct solvers add value through competition.} On CoW Protocol, trades routed through distinct solver addresses (not the settlement contract) achieve +3.57 bps average markout versus +3.07 bps for 1inch settlement trades (t=4.42, p$<$0.001). This +0.50 bps difference represents the combined benefit of batch auction mechanism design and visible solver competition. The competition effect is largest in falling markets (+1.18 bps), precisely when execution is most challenging.

    \item \textbf{Momentum dominates markout distribution.} We find strong correlation (0.54--0.62) between backward and forward markouts for major token pairs. Prices that were rising before execution tend to continue rising after, generating positive markouts. This momentum effect, not adverse selection or solver timing, is the primary determinant of execution quality.

    \item \textbf{Higher volatility correlates with better markouts.} Contrary to Loss-Versus-Rebalancing (LVR) theory \citep{milionis2022}, which predicts volatility should hurt execution quality, we find markouts increase from +1.27 bps in the lowest volatility quintile to +7.67 bps in the highest. We resolve this apparent contradiction by showing that high-volatility periods in our sample were predominantly uptrending markets, and 79\% of trades were sells that benefit from rising prices.

    \item \textbf{CEX-DEX latency arbitrage timing fails.} Testing the hypothesis that solvers observe CEX prices and time fills favorably, we find the opposite: ``favorable'' timing (e.g., buying when price just fell) produces \textit{worse} markouts because momentum continues. Solvers cannot profitably exploit short-term mean reversion.
\end{enumerate}

Our findings contribute to multiple literatures. For OFA design, we provide evidence that batch auctions and solver competition both matter, but momentum dynamics dominate execution outcomes. For market microstructure theory, we show that LVR predictions must be conditioned on market regime---trending markets can produce positive markouts despite high volatility. For practitioners, we demonstrate that solver timing strategies based on CEX price observation are not profitable.

The remainder of this paper is organized as follows. Section 2 reviews the relevant literature on OFAs, market microstructure, and MEV. Section 3 describes our data and methodology. Section 4 presents our main empirical results. Section 5 provides theoretical interpretation and connects findings to existing models. Section 6 discusses implications and concludes.

%==============================================================================
\section{Literature Review}
%==============================================================================

Our research connects to three strands of literature: order flow auctions and MEV in DeFi, traditional market microstructure, and execution quality measurement.

\subsection{Order Flow Auctions and MEV}

The concept of Maximal Extractable Value (MEV) was introduced by \cite{daian2020}, who documented how miners could extract value by reordering, inserting, or censoring transactions. This ``Flash Boys 2.0'' phenomenon created significant costs for DeFi users through frontrunning and sandwich attacks.

Order Flow Auctions emerged as a response to MEV extraction. \cite{lu2024} provides a comprehensive mapping of Ethereum's order flow landscape, identifying a four-layer sequential auction model: (1) solver auctions, where specialized agents compete to fill orders; (2) market maker auctions via Request-for-Quote (RFQ) systems; (3) order flow auctions that match users with searchers; and (4) block builder auctions that aggregate transactions. Lu documents 22 active solvers accessing 33+ liquidity sources, with profit margins ``clustering towards zero'' due to competition.

\cite{flashbots2023ofa1} warn about centralization risks from exclusive order flow (EOF) arrangements, identifying a feedback loop where wallets preferentially route to builders with highest inclusion rates. \cite{flashbots2023ofa2} analyze two fundamental OFA designs: explicit auctions where users withhold execution permission while extractors bid, and fee escalators that create implicit Dutch auctions through dynamically adjusting fees. Both designs suffer from incomplete information, creating tradeoffs between latency and value capture.

\cite{ma2025} provide game-theoretic analysis of OFAs under Proposer-Builder Separation (PBS), proving existence and uniqueness of Nash equilibrium in a two-player game between solvers and builders. Their key prediction---that builder centralization emerges endogenously---aligns with our empirical finding of a shared solver operating across platforms.

\cite{resnick2023} analyzes contingent fees in OFAs, while \cite{roughgarden2024} examine centralization in block building more broadly. The Frontier Research taxonomy \citep{frontier2024} identifies four distinct OFA designs based on purpose (price discovery vs. execution) and order type (swaps vs. all transactions), predicting market clustering around specialized implementations.

\subsection{Solver Concentration and Market Structure}

Recent research has documented severe concentration in solver and filler markets. \cite{flashbots2023ofa1} warn that ``two builders hold over 90\% of the block market'' and note that in November 2023, ``90\% of volume on UniswapX was filled by four incumbent market makers, with Wintermute filling 60\%.'' They identify a natural centralizing force: ``market making has historically been an aggressive centralizing force in both TradFi and DeFi due to economies of scale and the transferability of edge.''

Industry reports confirm these patterns. Blockworks \citep{blockworks2025} documents that Barter, the largest CoW Protocol solver, holds approximately 28\% market share and recently acquired competitor Copium Capital's codebase, further consolidating the market. The Sei Network \citep{sei2025} describes a ``centralization paradox'' where ``while solvers democratize access, economic forces may concentrate power among the most sophisticated players,'' creating ``winner-take-most dynamics.''

Notably, 1inch Fusion's resolver system has a \textbf{hard cap of 10 resolvers} by design, requiring substantial staking of 1INCH tokens (minimum 5\% of total Unicorn Power supply) and KYC/KYB verification. This structural constraint creates conditions for concentration that we empirically document.

Our contribution to this literature is the first systematic empirical analysis of solver concentration across platforms, quantifying the welfare cost of monopoly market structure.

\subsection{Traditional Market Microstructure}

Our analysis draws on foundational market microstructure theory. \cite{kyle1985} models how informed trading affects prices, predicting that order flow conveys information and moves prices permanently. We test Kyle's predictions by examining whether large orders show greater price impact, finding mixed evidence---BUY whales show potential informed trading against them (44.7\% correct direction), but overall, momentum rather than information drives our markout patterns.

\cite{cohen2012} develop a limit order book model for latency arbitrage, providing theoretical grounding for CEX-DEX arbitrage strategies. Their framework predicts that faster traders can exploit stale prices. We test this directly by examining whether solvers time fills based on CEX price movements, finding that momentum dominates any potential timing gains.

\cite{menkveld2016} surveys high-frequency trading economics, documenting how speed advantages affect market quality. While OFA solvers operate at slower timescales than traditional HFT, similar dynamics of information asymmetry and latency apply.

\subsection{Loss-Versus-Rebalancing (LVR)}

\cite{milionis2022} introduce the Loss-Versus-Rebalancing framework for AMMs, showing that liquidity providers systematically lose to arbitrageurs who correct stale prices. LVR increases with volatility---a prediction we directly test. Our finding that higher volatility correlates with \textit{better} markouts appears to contradict LVR theory, but we show this reflects sample composition (trending markets with sell-heavy flow) rather than a fundamental theoretical failure.

\subsection{Execution Quality Measurement}

\cite{bachu2024} propose methodology for quantifying price improvement in OFAs, decomposing improvements into routing efficiency, gas optimization, and other factors. They document 4--5 basis points improvement over Uniswap router for 1inch and UniswapX.

Most directly relevant is \cite{yuminaga2025}, who analyze execution welfare across solver-based DEXes using data from August 2023 to February 2024. They find: (1) negative markouts relative to Binance across all platforms, suggesting onchain markets have lower price efficiency; (2) negligible correlation (0.03--0.07) between volatility and execution welfare; and (3) CoWSwap outperforms at scale, particularly for large trades.

Our study differs from Yuminaga et al. in several critical ways, summarized in Table \ref{tab:yuminaga_comparison}.

\begin{table}[H]
\centering
\caption{Comparison with Yuminaga et al. (2025)}
\label{tab:yuminaga_comparison}
\begin{tabular}{lll}
\toprule
\textbf{Dimension} & \textbf{Yuminaga et al.} & \textbf{Our Study} \\
\midrule
Sample Period & Aug 2023 -- Feb 2024 & Jun -- Sep 2025 \\
Market Regime & Bear market recovery & Bull market \\
Avg Markout vs Binance & Negative (-1000 to -5 bps) & Positive (+3.3 bps) \\
Volatility Correlation & Near zero (0.03--0.07) & Strong (0.54--0.62) \\
Outlier Handling & No IQR filter & IQR filter [-24, +33] bps \\
Order Flow Composition & Not reported & 79\% SELLs \\
\bottomrule
\end{tabular}
\end{table}

We attribute the sign reversal in markouts to four factors: (1) \textbf{Market regime}: our sample captures a bull market where SELLs benefit from rising prices; (2) \textbf{Order flow composition}: 79\% of our trades are SELLs, which generate positive markouts in rising markets; (3) \textbf{IQR filtering}: we remove extreme outliers that Yuminaga includes (their -1000 bps observations); (4) \textbf{Sample period}: 18+ months of solver infrastructure evolution may have improved execution quality.

The volatility correlation discrepancy is explained by sample composition: Yuminaga's low correlation reflects mixed market conditions, while our strong correlation reflects predominantly trending (momentum-driven) markets where volatility proxies for trend strength.

\cite{uniswap2024} establish methodology for comparing OFAs, finding UniswapX provides approximately 4.7 bps average price improvement. The CAAW 2025 paper extends this analysis with additional decomposition frameworks.

\subsection{Batch Auction Mechanism Design}

Our finding that CoW Protocol outperforms 1inch Fusion by +0.38 bps even for trades using the same settlement infrastructure is consistent with theoretical predictions about batch auction advantages. CoW Protocol's batch auction mechanism provides three key benefits over 1inch Fusion's Dutch auction:

\begin{enumerate}
    \item \textbf{Uniform Clearing Prices}: All orders in a batch trade at the same price, eliminating transaction ordering games that enable MEV extraction \citep{cow2025}.
    \item \textbf{Coincidence of Wants (CoW) Matching}: Direct peer-to-peer matching reduces reliance on external liquidity and associated slippage.
    \item \textbf{Multi-Asset Batching}: Unlike Dutch auctions that handle one token at a time, batch auctions can optimize across multiple assets simultaneously.
\end{enumerate}

The July 2025 introduction of Fair Combinatorial Batch Auctions (FCBA) \citep{canidio2024} further improved CoW Protocol's mechanism by allowing multiple solvers to win simultaneously while guaranteeing each trader receives at least what they would in simultaneous first-price auctions.

\subsection{Contribution to Literature}

Our paper makes several contributions:

\begin{enumerate}[leftmargin=*]
    \item \textbf{Data artifact identification}: We discover that 1inch Fusion trades in our data route through the GPv2Settlement contract, preventing resolver-level analysis. This finding is important for future researchers using similar data sources.

    \item \textbf{Competition value quantification}: We quantify the value of solver competition at +0.50 bps per trade (t=4.42), comparing distinct CoW solvers to settlement-level trades. Competition matters most in falling markets (+1.18 bps advantage).

    \item \textbf{Three-way comparison methodology}: We introduce a three-way comparison (CoW Distinct Solvers vs CoW Settlement vs 1inch Settlement) that decomposes the CoW advantage into mechanism design (+0.38 bps) and solver-level competition (+0.12 bps) components.

    \item \textbf{Momentum identification}: We are first to document that price momentum, not adverse selection, is the primary driver of OFA markout distribution.

    \item \textbf{LVR reconciliation}: We explain why volatility can correlate positively with markouts in trending markets, reconciling apparent contradictions with theory.

    \item \textbf{CEX-DEX timing test}: We provide first explicit test of whether solvers exploit CEX price movements, finding they cannot.

    \item \textbf{Market-state conditioning}: We show that competition effects are heterogeneous across market conditions, with the largest benefits during falling markets where execution is most challenging.

    \item \textbf{Backward markout methodology}: We introduce analysis of pre-trade price dynamics to characterize positive vs. negative markout trades.
\end{enumerate}

%==============================================================================
\section{Data and Methodology}
%==============================================================================

\subsection{Data Sources}

Our analysis combines data from three sources:

\textbf{On-chain execution data}: We collect all transactions routed through CoW Protocol and 1inch Fusion during four sampling periods:
\begin{itemize}
    \item Period 1 (P1): June 11--18, 2025
    \item Period 2 (P2): July 11--18, 2025
    \item Period 3 (P3): August 10--17, 2025
    \item Period 4 (P4): August 26 -- September 2, 2025
\end{itemize}

Each period spans 8 days, yielding 32 total trading days. We extract transaction hashes, timestamps, token addresses, amounts, and solver addresses from blockchain data.

\textbf{CEX price data}: We obtain millisecond-level mid-prices from Binance for all traded pairs. Binance serves as the price benchmark due to its dominant market share and tight spreads for major cryptocurrency pairs.

\textbf{Token mapping}: We map on-chain token addresses to Binance trading symbols, successfully matching 18 token pairs representing the majority of OFA volume.

\subsection{Sample Construction}

Our initial dataset contains 82,657 transactions. We apply the following filters:

\begin{enumerate}[leftmargin=*]
    \item \textbf{Binance match}: Retain only trades with corresponding Binance price data (removes long-tail tokens without CEX listings).

    \item \textbf{IQR outlier filter}: Remove trades with 60-second markouts outside the interquartile range bounds of $[-24, +33]$ bps. This removes 15.8\% of observations, yielding 69,661 trades for analysis.

    \item \textbf{Native ETH exclusion}: We identify that native ETH trades (address 0xeee...eee) are 100\% buy-directional due to a data artifact, and account for this in directional analysis.
\end{enumerate}

Table \ref{tab:sample_summary} presents sample composition.

\begin{table}[H]
\centering
\caption{Sample Summary Statistics}
\label{tab:sample_summary}
\begin{tabular}{lrr}
\toprule
\textbf{Metric} & \textbf{CoW Protocol} & \textbf{1inch Fusion} \\
\midrule
Total Trades & 39,009 & 30,652 \\
Period 1 (P1) & 6,934 & 5,160 \\
Period 2 (P2) & 7,336 & 5,821 \\
Period 3 (P3) & 12,746 & 9,830 \\
Period 4 (P4) & 11,993 & 9,841 \\
\midrule
Avg Markout (bps) & +3.59 & +2.95 \\
Pct Positive & 64.5\% & 63.4\% \\
Pct SELL & 78.1\% & 80.7\% \\
\bottomrule
\end{tabular}
\end{table}

\subsection{Markout Methodology}

We employ forward and backward markouts to measure execution quality and price dynamics around trade execution. For a trade executed at time $t$ with execution price $P_t^{exec}$:

\textbf{Forward markout} measures subsequent price movement:
\begin{equation}
M_t^{+\tau} = \frac{P_{t+\tau}^{CEX} - P_t^{exec}}{P_t^{exec}} \times 10000 \text{ (in bps)}
\end{equation}

We compute forward markouts at $\tau \in \{1s, 10s, 60s\}$. A positive markout for a BUY order indicates the price rose after execution---the user bought before a price increase, representing good execution. For SELL orders, a positive markout indicates the price rose after selling, which also represents good execution (sold before price increased).

\textbf{Backward markout} measures prior price movement:
\begin{equation}
M_t^{-\tau} = \frac{P_t^{exec} - P_{t-\tau}^{CEX}}{P_{t-\tau}^{CEX}} \times 10000 \text{ (in bps)}
\end{equation}

Backward markouts capture price dynamics leading into order execution, allowing us to analyze momentum and mean reversion.

\textbf{Data quality note}: We discovered that markout\_1s = markout\_10s for 99.9\% of observations, indicating these fields contain identical values. Meaningful variation exists only between the 10s and 60s windows, so we focus analysis on 60-second markouts.

\subsection{Direction Classification}

We classify trade direction based on the output token:
\begin{itemize}
    \item \textbf{BUY}: Output token is ETH (native or WETH)
    \item \textbf{SELL}: Output token is stablecoin or other token
\end{itemize}

This classification reflects the user's intent relative to ETH exposure. Approximately 79\% of trades are classified as SELLs, indicating users predominantly exit ETH positions through OFAs.

\subsection{Volatility and Market State Classification}

We construct several market condition variables:

\textbf{Volatility proxy}: $|M_t^{-60s}|$, the absolute backward markout, captures recent price variation.

\textbf{Market state}: Based on backward markout:
\begin{itemize}
    \item RISING: $M_t^{-60s} > 5$ bps
    \item FALLING: $M_t^{-60s} < -5$ bps
    \item STABLE: $|M_t^{-60s}| \leq 5$ bps
\end{itemize}

\textbf{Volatility quintiles}: We sort trades into quintiles based on $|M_t^{-60s}|$ to test the volatility-execution quality relationship.

\subsection{Platform and Period Fixed Effects}

Our panel structure allows us to estimate specifications with platform and period fixed effects:
\begin{equation}
M_{i,t}^{+60s} = \alpha + \beta \cdot \text{CoW}_i + \gamma \cdot X_{i,t} + \delta_p + \epsilon_{i,t}
\end{equation}

where $\text{CoW}_i$ indicates CoW Protocol, $X_{i,t}$ are controls, and $\delta_p$ are period fixed effects. This addresses concerns about time-varying market conditions confounding platform comparisons.

%==============================================================================
\section{Empirical Results}
%==============================================================================

\subsection{Overview: Markout Distribution}

Table \ref{tab:main_markouts} presents main markout statistics by platform and period.

\begin{table}[H]
\centering
\caption{Forward 60-Second Markouts by Platform and Period}
\label{tab:main_markouts}
\begin{tabular}{llrrrrr}
\toprule
\textbf{Period} & \textbf{Platform} & \textbf{N} & \textbf{Mean} & \textbf{Std} & \textbf{Pct Pos} & \textbf{Momentum Corr} \\
\midrule
P1 & CoW & 6,934 & +3.67 & 10.49 & 65.7\% & 0.532 \\
P1 & 1inch & 5,160 & +3.20 & 10.01 & 65.3\% & 0.589 \\
P2 & CoW & 7,336 & +4.20 & 11.56 & 64.7\% & 0.042 \\
P2 & 1inch & 5,821 & +3.84 & 11.26 & 64.1\% & 0.042 \\
P3 & CoW & 12,746 & +4.09 & 10.50 & 65.4\% & 0.481 \\
P3 & 1inch & 9,830 & +3.90 & 10.22 & 65.2\% & 0.526 \\
P4 & CoW & 11,993 & +2.30 & 9.07 & 62.4\% & 0.074 \\
P4 & 1inch & 9,841 & +1.73 & 8.66 & 60.7\% & 0.073 \\
\midrule
\textbf{Total} & \textbf{CoW} & \textbf{39,009} & \textbf{+3.59} & \textbf{10.28} & \textbf{64.5\%} & \textbf{---} \\
\textbf{Total} & \textbf{1inch} & \textbf{30,652} & \textbf{+2.95} & \textbf{9.88} & \textbf{63.4\%} & \textbf{---} \\
\bottomrule
\end{tabular}
\end{table}

Key observations:
\begin{itemize}
    \item Average markouts are \textbf{positive} across all periods and platforms (+2.95 to +3.59 bps)
    \item CoW Protocol consistently outperforms 1inch Fusion by 0.36--0.64 bps
    \item Periods 1 and 3 show high momentum correlation (0.48--0.59), while P2 and P4 show near-zero correlation
    \item Positive markout rate stable at 60--66\%
\end{itemize}

\subsection{Finding 1: Momentum Dominates Markout Distribution}

We test whether backward price movements predict forward markouts. Table \ref{tab:momentum} presents results by market state.

\begin{table}[H]
\centering
\caption{Forward Markouts by Prior Market State}
\label{tab:momentum}
\begin{tabular}{lrrrrr}
\toprule
\textbf{Prior Move} & \textbf{N} & \textbf{Avg Bwd 60s} & \textbf{Avg Fwd 60s} & \textbf{Pct Pos} & \textbf{Same Sign \%} \\
\midrule
EXTREME\_FALL & 79 & -2,303 & -1.17 & 39.2\% & 60.8\% \\
LARGE\_FALL & 1,457 & -27.5 & -6.38 & 30.5\% & 69.5\% \\
SMALL\_FALL & 8,964 & -10.3 & -4.82 & 27.5\% & 72.5\% \\
STABLE & 36,524 & +0.54 & +1.61 & 61.2\% & 70.7\% \\
SMALL\_RISE & 18,929 & +10.5 & +8.16 & 84.1\% & 84.1\% \\
LARGE\_RISE & 3,583 & +27.6 & +18.9 & 91.5\% & 91.5\% \\
EXTREME\_RISE & 125 & +216 & +13.0 & 71.2\% & 54.4\% \\
\bottomrule
\end{tabular}
\end{table}

\textbf{Key finding}: Price movements persist strongly. When the prior 60 seconds showed a rise, prices continue rising 84--92\% of the time. Mean reversion only appears at extremes (EXTREME\_FALL and EXTREME\_RISE show negative correlation with forward moves).

Token-level momentum correlations reveal substantial heterogeneity:

\begin{table}[H]
\centering
\caption{Momentum Correlation by Token Type}
\label{tab:token_momentum}
\begin{tabular}{llrrrr}
\toprule
\textbf{Token Type} & \textbf{Platform} & \textbf{N} & \textbf{Avg Markout} & \textbf{Pct Pos} & \textbf{Momentum Corr} \\
\midrule
SHORT\_TAIL & CoW & 32,897 & +3.78 & 66.5\% & 0.538 \\
SHORT\_TAIL & 1inch & 25,097 & +3.35 & 65.8\% & 0.621 \\
LONG\_TAIL & CoW & 4,595 & +2.63 & 56.2\% & 0.032 \\
LONG\_TAIL & 1inch & 4,506 & +2.66 & 56.3\% & 0.032 \\
MAJOR\_CROSS & CoW & 1,517 & -0.31 & 44.9\% & 0.634 \\
MAJOR\_CROSS & 1inch & 1,049 & -1.82 & 39.8\% & 0.620 \\
\bottomrule
\end{tabular}
\end{table}

Short-tail tokens (USDC/USDT-ETH/BTC pairs) show strong momentum (0.54--0.62), while long-tail altcoins show essentially zero momentum correlation. Major crosses (ETH/BTC) show the strongest momentum but negative average markouts, reflecting that these pairs had predominantly falling prices during our sample.

\subsection{Finding 2: Volatility Correlates with Better Markouts}

Table \ref{tab:volatility} tests the relationship between volatility and execution quality.

\begin{table}[H]
\centering
\caption{Markouts by Volatility Quintile}
\label{tab:volatility}
\begin{tabular}{lrrrrrrr}
\toprule
\textbf{Quintile} & \textbf{N} & \textbf{Avg Vol} & \textbf{Markout} & \textbf{Pct Rising} & \textbf{If Rising} & \textbf{If Falling} & \textbf{Pct Sells} \\
\midrule
Q1 (lowest) & 13,933 & 0.51 & +1.27 & 51.0\% & +1.57 & +0.96 & 87.7\% \\
Q2 & 13,932 & 2.16 & +1.61 & 59.2\% & +2.44 & +0.39 & 77.7\% \\
Q3 & 13,932 & 4.69 & +2.27 & 67.7\% & +4.19 & -1.75 & 76.6\% \\
Q4 & 13,932 & 9.08 & +3.71 & 68.1\% & +7.29 & -3.91 & 78.7\% \\
Q5 (highest) & 13,932 & 34.16 & +7.67 & 69.0\% & +14.33 & -7.19 & 75.6\% \\
\bottomrule
\end{tabular}
\end{table}

\textbf{Apparent contradiction with LVR}: Higher volatility should increase adverse selection costs per \cite{milionis2022}, producing worse markouts. We observe the opposite: Q5 markouts (+7.67 bps) are 6x higher than Q1 (+1.27 bps).

\textbf{Resolution}: The correlation operates through market regime:
\begin{enumerate}
    \item High-volatility periods had more rising markets (69\% vs 51\%)
    \item Rising markets produce positive markouts due to momentum
    \item 79\% of trades are SELLs, which benefit from rising prices
\end{enumerate}

Within rising markets, higher volatility still produces better markouts (+14.33 vs +1.57 bps). But within falling markets, higher volatility produces \textit{worse} markouts (-7.19 vs +0.96 bps), consistent with LVR. The overall positive correlation reflects that our sample was predominantly uptrending.

\subsection{Finding 3: Settlement Contract vs. Solver-Level Analysis}

Our analysis of solver addresses reveals a critical data artifact that informs how we interpret platform differences.

\begin{table}[H]
\centering
\caption{Address Distribution by Platform}
\label{tab:addresses}
\begin{tabular}{llrrr}
\toprule
\textbf{Platform} & \textbf{Address} & \textbf{Trades} & \textbf{Share} & \textbf{Avg Markout} \\
\midrule
\textbf{1inch Fusion} & 0x9008D19f58... & 30,652 & 100.0\% & +3.07 bps \\
\midrule
CoW Protocol & 0x9008D19f58... & 26,400 & 67.68\% & +3.45 bps \\
CoW Protocol & 0x4dd1be0Cd6... & 5,420 & 13.89\% & +4.03 bps \\
CoW Protocol & 0xA9D635EF85... & 5,158 & 13.22\% & +3.76 bps \\
CoW Protocol & 0x00806DaA2C... & 1,543 & 3.96\% & +1.41 bps \\
CoW Protocol & 0x6bf97aFe2D... & 265 & 0.68\% & +3.93 bps \\
CoW Protocol & 0x95480d3f27... & 223 & 0.57\% & +2.46 bps \\
\bottomrule
\end{tabular}
\end{table}

\textbf{Critical discovery}: The address 0x9008D19f58... is the \textbf{GPv2Settlement Contract}---the smart contract through which all CoW Protocol trades are settled. This address appearing on 100\% of 1inch Fusion trades indicates a \textit{data artifact}, not a monopoly:

\begin{itemize}
    \item \textbf{1inch Fusion}: Our data captures settlement-level information, not resolver identity. Industry reports (Messari Q2 2025) indicate 1inch has multiple active resolvers with the top 3 controlling approximately 68\% of volume.
    \item \textbf{CoW Protocol}: Our data captures both settlement contract trades (67.68\%) and distinct solver addresses (32.32\%), enabling solver-level analysis.
\end{itemize}

\textbf{Revised interpretation}: Rather than claiming a monopoly, we can construct a three-way comparison:

\begin{table}[H]
\centering
\caption{Three-Way Execution Quality Comparison}
\label{tab:threeway}
\begin{tabular}{lrrr}
\toprule
\textbf{Trade Type} & \textbf{Trades} & \textbf{Avg Markout} & \textbf{Pct Positive} \\
\midrule
CoW Distinct Solvers & 12,609 & \textbf{+3.57 bps} & 64.5\% \\
CoW Settlement & 26,400 & +3.45 bps & 64.4\% \\
1inch (Settlement) & 30,652 & +3.07 bps & 63.5\% \\
\bottomrule
\end{tabular}
\end{table}

\textbf{Implications}:
\begin{enumerate}
    \item \textbf{Mechanism design matters}: CoW Settlement outperforms 1inch Settlement by +0.38 bps, suggesting batch auction design provides value independent of solver identity.
    \item \textbf{Solver-level competition adds value}: CoW Distinct Solvers outperform CoW Settlement by +0.12 bps, isolating the competition effect.
    \item \textbf{Combined benefit}: CoW Distinct Solvers outperform 1inch by +0.50 bps---the total advantage of mechanism design plus visible competition.
\end{enumerate}

\subsection{Finding 4: Competition Adds Measurable Value}

We test whether the execution quality differences between trade types are statistically significant using a two-sample t-test.

\begin{table}[H]
\centering
\caption{Statistical Comparison: CoW Distinct Solvers vs 1inch Settlement}
\label{tab:ttest}
\begin{tabular}{lrrrrrrr}
\toprule
\textbf{Group 1} & \textbf{Group 2} & \textbf{N$_1$} & \textbf{N$_2$} & \textbf{Mean$_1$} & \textbf{Mean$_2$} & \textbf{Diff} & \textbf{t-stat} \\
\midrule
CoW Distinct & 1inch & 12,609 & 30,652 & +3.57 & +3.07 & \textbf{+0.49} & \textbf{4.42} \\
\bottomrule
\end{tabular}
\end{table}

\textbf{Key result}: CoW Distinct Solvers provide \textbf{+0.49 bps better execution} than 1inch Settlement trades (t=4.42, p$<$0.001). This is \textbf{highly statistically significant}.

\textbf{All three trade types significantly different from zero}:

\begin{table}[H]
\centering
\caption{T-Tests vs Zero for Each Trade Type}
\label{tab:ttest_zero}
\begin{tabular}{lrrrrr}
\toprule
\textbf{Trade Type} & \textbf{N} & \textbf{Mean (bps)} & \textbf{Std Dev} & \textbf{t-stat vs 0} & \textbf{Significance} \\
\midrule
CoW Distinct Solvers & 12,609 & +3.57 & 10.79 & 37.12 & p$<$0.001 \\
CoW Settlement & 26,400 & +3.45 & 10.10 & 55.46 & p$<$0.001 \\
1inch Settlement & 30,652 & +3.07 & 9.97 & 53.97 & p$<$0.001 \\
\bottomrule
\end{tabular}
\end{table}

\textbf{Competition effect by market condition}:

\begin{table}[H]
\centering
\caption{Competition Effect Across Market States}
\label{tab:competition_by_market}
\begin{tabular}{lrrrr}
\toprule
\textbf{Market State} & \textbf{1inch} & \textbf{CoW Settlement} & \textbf{CoW Distinct} & \textbf{Competition Benefit} \\
\midrule
Falling & -5.67 bps & -4.50 bps & -4.49 bps & \textbf{+1.17 to +1.18 bps} \\
Stable & +1.40 bps & +1.51 bps & +2.45 bps & +0.11 to +1.05 bps \\
Rising & +9.98 bps & +9.60 bps & +10.34 bps & -0.38 to +0.36 bps \\
\bottomrule
\end{tabular}
\end{table}

\textbf{Critical finding}: \textbf{Competition matters most in falling markets} where execution is hardest. When prices are falling:
\begin{itemize}
    \item 1inch: -5.67 bps, 25.4\% positive trades
    \item CoW Distinct: -4.49 bps, 30.0\% positive trades
    \item \textbf{Benefit: +1.18 bps and +4.6 percentage points more positive trades}
\end{itemize}

This suggests that mechanism design and solver competition provide the greatest value precisely when users need it most---during adverse market conditions.

\subsection{Finding 5: CEX-DEX Timing Cannot Be Exploited}

We test whether solvers time fills based on CEX price movements. If solvers observe CEX prices and fill orders when moves are ``favorable'' (price fell before BUY, price rose before SELL), this should produce better markouts.

\begin{table}[H]
\centering
\caption{Markouts by CEX Move Direction and Order Type}
\label{tab:timing}
\begin{tabular}{llrrrl}
\toprule
\textbf{Direction} & \textbf{CEX Move} & \textbf{N} & \textbf{Markout} & \textbf{Pct Pos} & \textbf{Theory Timing} \\
\midrule
BUY & PRICE\_FELL & 5,483 & -8.97 & 35.2\% & FAVORABLE \\
BUY & PRICE\_ROSE & 8,956 & +11.70 & 87.1\% & UNFAVORABLE \\
SELL & PRICE\_FELL & 21,372 & -6.35 & 27.9\% & UNFAVORABLE \\
SELL & PRICE\_ROSE & 33,850 & +8.49 & 85.7\% & FAVORABLE \\
\bottomrule
\end{tabular}
\end{table}

\textbf{Counter to timing theory}: ``Favorable'' timing produces WORSE markouts for BUY orders (-8.97 vs +11.70). The reason is momentum: if price fell before a BUY, it continues falling after, generating negative markouts. Solvers cannot profitably time market reversals because short-term momentum dominates mean reversion.

All solvers show identical patterns regardless of platform:

\begin{table}[H]
\centering
\caption{Solver Performance by Market State (All Show Identical Patterns)}
\label{tab:solver_timing}
\begin{tabular}{llrr}
\toprule
\textbf{Solver} & \textbf{Market State} & \textbf{Avg Markout} & \textbf{Pct Positive} \\
\midrule
0x9008D19f... (shared) & FALLING & -5.57 & 24.6\% \\
0x9008D19f... (shared) & STABLE & +1.40 & 60.7\% \\
0x9008D19f... (shared) & RISING & +8.21 & 85.0\% \\
\midrule
0x4dd1be0C... (CoW-only) & FALLING & -1.48 & 44.2\% \\
0x4dd1be0C... (CoW-only) & STABLE & +2.77 & 67.4\% \\
0x4dd1be0C... (CoW-only) & RISING & +7.87 & 79.1\% \\
\bottomrule
\end{tabular}
\end{table}

No solver demonstrates ability to beat market conditions through timing.

\subsection{Additional Results: Trade Size, Hourly Patterns, and Distributional Effects}

\subsubsection{Trade Size Effects}

Literature predicts larger trades benefit more from solver routing \citep{yuminaga2025}. We find a different pattern:

\begin{table}[H]
\centering
\caption{Markouts by Trade Size}
\label{tab:trade_size}
\begin{tabular}{lrrrrr}
\toprule
\textbf{Size Bucket} & \textbf{N} & \textbf{CoW} & \textbf{1inch} & \textbf{Std Dev} & \textbf{Pct Pos} \\
\midrule
<\$100 & 2,692 & +0.54 & -0.04 & 13.10 & 49.2\% \\
\$100--1K & 8,064 & +6.00 & +5.69 & 12.38 & 68.9\% \\
\$1K--10K & 20,375 & +4.09 & +3.60 & 9.85 & 64.1\% \\
\$10K--50K & 16,609 & +3.26 & +3.39 & 8.63 & 65.9\% \\
\$50K--100K & 3,313 & +3.09 & +3.24 & 8.84 & 67.3\% \\
>\$100K & 12,449 & +3.67 & +3.77 & 7.64 & 70.6\% \\
\bottomrule
\end{tabular}
\end{table}

\textbf{Key findings}:
\begin{itemize}
    \item \$100--1K bracket has HIGHEST markouts (+6.00 bps), not largest trades
    \item Very small trades (<\$100) have near-zero markouts
    \item Variance decreases with size (large trades more predictable)
    \item CoW advantage disappears for large trades (platforms converge)
\end{itemize}

\subsubsection{Hourly Patterns}

Markouts are remarkably stable across trading hours:

\begin{table}[H]
\centering
\caption{Markouts by Trading Session}
\label{tab:hourly}
\begin{tabular}{lrrr}
\toprule
\textbf{Session} & \textbf{Hours (UTC)} & \textbf{Avg Markout} & \textbf{Pct Positive} \\
\midrule
Asia & 22:00--07:00 & +3.21 & 62.7\% \\
Europe & 08:00--16:00 & +3.29 & 65.6\% \\
US & 17:00--21:00 & +3.38 & 63.4\% \\
\bottomrule
\end{tabular}
\end{table}

Europe session shows slightly higher positive markout rate (65.6\%), but differences are economically small.

\subsubsection{Who Benefits from Positive Markouts?}

We characterize the distributional effects by comparing positive and negative markout trades:

\begin{table}[H]
\centering
\caption{Characteristics of Positive vs Negative Markout Trades}
\label{tab:distribution}
\begin{tabular}{llrrrr}
\toprule
\textbf{Markout Sign} & \textbf{Direction} & \textbf{N} & \textbf{Avg Fwd} & \textbf{Avg Bwd} & \textbf{\% of Dir} \\
\midrule
POSITIVE & BUY & 8,951 & +10.70 & +5.75 & 62.0\% \\
POSITIVE & SELL & 35,635 & +8.15 & +4.53 & 64.5\% \\
NEGATIVE & BUY & 5,488 & -7.70 & -12.64 & 38.0\% \\
NEGATIVE & SELL & 19,587 & -5.88 & -6.55 & 35.5\% \\
\bottomrule
\end{tabular}
\end{table}

\textbf{Key insight}: Positive markout trades had prices RISING before execution (+4.5 to +5.8 bps backward). Negative markout trades had prices FALLING before (-6.6 to -12.6 bps backward). BUY orders during downtrends (-12.64 bps prior) experience especially poor outcomes.

\subsection{Robustness: Cross-Period Stability}

Our findings are stable across periods, though momentum correlation varies with market regime:

\begin{itemize}
    \item \textbf{P1 and P3} (high momentum correlation 0.48--0.59): Trending markets
    \item \textbf{P2 and P4} (low momentum correlation 0.04--0.07): Range-bound/choppy markets
    \item CoW advantage persists in all periods
    \item Positive markout rate stable at 60--66\%
\end{itemize}

The variation in momentum correlation suggests our findings on momentum dominance are strongest during trending periods. In choppy markets, neither momentum nor timing strategies would be effective.

\subsection{Finding 6: Token Pair Directional Asymmetry}

We discover a striking directional asymmetry in execution quality across token pairs. This analysis examines how trade direction---specifically whether users are selling ETH for stablecoins or buying ETH with stablecoins---affects markout outcomes.

\begin{table}[H]
\centering
\caption{Token Pair Performance by Direction}
\label{tab:token_pair}
\begin{tabular}{lrrrl}
\toprule
\textbf{Token Pair} & \textbf{Markout (\%)} & \textbf{Win Rate} & \textbf{N} & \textbf{Direction} \\
\midrule
WETH$\rightarrow$USDT & +0.252 & 75.5\% & 8,234 & ETH $\rightarrow$ Stable \\
WETH$\rightarrow$USDC & +0.235 & 78.6\% & 12,188 & ETH $\rightarrow$ Stable \\
USDC$\rightarrow$USDT & +0.011 & 55.4\% & 15,498 & Stable $\leftrightarrow$ Stable \\
USDT$\rightarrow$USDC & -0.011 & 46.6\% & 15,491 & Stable $\leftrightarrow$ Stable \\
USDC$\rightarrow$WETH & -0.312 & 26.6\% & 5,708 & Stable $\rightarrow$ ETH \\
USDT$\rightarrow$WETH & -0.484 & 27.9\% & 3,238 & Stable $\rightarrow$ ETH \\
\bottomrule
\end{tabular}
\end{table}

\textbf{Key findings}:

\begin{enumerate}[leftmargin=*]
    \item \textbf{Directional asymmetry is dramatic}: ETH$\rightarrow$Stablecoin trades achieve +0.24\% average markout, while Stablecoin$\rightarrow$ETH trades show -0.40\% average markout---a \textbf{64 basis point spread}.

    \item \textbf{Win rates diverge by 50 percentage points}: ETH sellers achieve 77\% positive markout rate versus only 27\% for ETH buyers. This is far larger than the platform effect (+1.1 pp) or solver competition effect (+4.6 pp in falling markets).

    \item \textbf{Stablecoin-to-stablecoin pairs are neutral}: USDC$\leftrightarrow$USDT trades show near-zero markouts ($\pm$0.01\%) and approximately 50\% win rates, as expected for stable pair arbitrage.
\end{enumerate}

\begin{figure}[H]
\centering
\includegraphics[width=\textwidth]{figures/figure4_token_pair_analysis.png}
\caption{Token Pair Performance Analysis. Panel (A) shows average markout by token pair direction, revealing strong directional asymmetry. Panel (B) shows win rates (probability of positive markout), with ETH$\rightarrow$Stablecoin trades achieving 77\% win rate versus 27\% for Stablecoin$\rightarrow$ETH trades.}
\label{fig:token_pair}
\end{figure}

\textbf{Interpretation}: This directional asymmetry strongly supports our momentum hypothesis. During our sample period (Summer 2025), markets were predominantly rising. Users \textit{selling} ETH for stablecoins were trading \textit{with} the trend---they sold before prices continued rising, generating positive markouts. Users \textit{buying} ETH were trading \textit{against} the trend---they bought before prices continued rising (which sounds favorable), but the momentum effect means prices had already risen before execution, locking in unfavorable entry prices.

\textbf{Implications for users}:
\begin{itemize}
    \item Trade direction relative to market trend is the \textbf{dominant factor} in execution quality
    \item Users selling into strength (ETH$\rightarrow$Stablecoin in rising markets) achieve consistently positive outcomes
    \item Users buying into strength (Stablecoin$\rightarrow$ETH in rising markets) face adverse momentum
    \item Platform choice (+0.50 bps) and solver competition (+0.12 bps) provide marginal improvements compared to the 64 bps directional effect
\end{itemize}

This finding reinforces our conclusion that market conditions dominate mechanism design in determining execution quality. The 64 bps directional spread dwarfs the 0.50 bps platform effect, suggesting that users should focus primarily on trade timing and direction rather than platform selection.

%==============================================================================
\section{Extended Analysis: Beyond Markouts}
%==============================================================================

To provide a more complete picture of OFA execution quality, we conduct six additional analyses that go beyond traditional markout measurement. These analyses examine surplus capture efficiency, execution consistency, order flow toxicity, solver specialization, cross-token correlations, and market regime effects.

\subsection{Surplus Capture Analysis}

We analyze how much of the ``available surplus''---defined as the absolute value of backward price movement---is captured by users in their forward markouts. This measures OFA efficiency in passing value to users.

\begin{table}[H]
\centering
\caption{Surplus Capture by Trade Pattern}
\label{tab:surplus}
\begin{tabular}{lrrrp{3cm}}
\toprule
\textbf{Pattern} & \textbf{N} & \textbf{Fwd Markout} & \textbf{Volatility} & \textbf{Interpretation} \\
\midrule
CONTINUED\_UP + SELL & 29,429 & +8.38 & 7.93 & Momentum continuation \\
CONTINUED\_DOWN + SELL & 13,919 & -6.33 & 12.10 & Momentum continuation \\
CONTINUED\_UP + BUY & 6,312 & +11.70 & 10.64 & Momentum continuation \\
REVERSED\_UP + SELL & 6,234 & +7.31 & 11.55 & Mean reversion \\
REVERSED\_DOWN + SELL & 5,640 & -4.82 & 8.10 & Mean reversion \\
CONTINUED\_DOWN + BUY & 2,990 & -8.98 & 28.41 & Adverse selection \\
\bottomrule
\end{tabular}
\end{table}

\textbf{Key findings}:
\begin{itemize}
    \item \textbf{Surplus Capture Ratio}: 34.7\% (users capture approximately one-third of available price movement)
    \item \textbf{Average Available Surplus}: $\sim$10.0 bps (absolute backward markout)
    \item \textbf{Average Captured Surplus}: $\sim$3.5 bps (forward markout)
\end{itemize}

The remaining $\sim$65\% of potential surplus is absorbed by solver margins, gas costs, execution slippage, and adverse selection in some cases. The CONTINUED\_DOWN + BUY pattern shows the highest volatility (28.4 bps), indicating these may be informed traders buying against the trend.

\subsection{Execution Consistency Analysis}

We measure the variability of execution quality using the coefficient of variation (CV) and distribution percentiles.

\begin{table}[H]
\centering
\caption{Execution Consistency Metrics by Platform}
\label{tab:consistency}
\begin{tabular}{lrrrrr}
\toprule
\textbf{Platform} & \textbf{Mean} & \textbf{Std Dev} & \textbf{CV} & \textbf{5th \%ile} & \textbf{95th \%ile} \\
\midrule
CoW Protocol & +3.59 & 10.28 & 2.86 & -13.3 & +20.5 \\
1inch Fusion & +2.95 & 9.88 & 3.35 & -13.3 & +19.2 \\
\bottomrule
\end{tabular}
\end{table}

\textbf{Key findings}:
\begin{itemize}
    \item \textbf{High uncertainty}: Coefficient of Variation $>$ 2 indicates substantial execution uncertainty
    \item \textbf{Negative tail risk}: 5th percentile at -13.3 bps means approximately 16\% of trades have negative markouts
    \item \textbf{Wide range}: 33--34 bps spread between 5th and 95th percentiles
    \item \textbf{CoW advantage}: Slightly lower CV (2.86 vs 3.35) indicates more consistent execution
\end{itemize}

The high variance means that while average markouts are positive, individual trade outcomes remain unpredictable. Users should not expect consistent positive markouts on every trade.

\subsection{Order Flow Toxicity Analysis}

We classify order flow into ``benign'' and ``potentially toxic'' categories based on the relationship between price direction and trade direction.

\begin{table}[H]
\centering
\caption{Order Flow Toxicity Classification}
\label{tab:toxicity}
\begin{tabular}{lrrrl}
\toprule
\textbf{Pattern} & \textbf{N} & \textbf{Markout} & \textbf{Classification} \\
\midrule
CONTINUED\_UP (SELL) & 29,429 & +8.38 & Benign \\
CONTINUED\_UP (BUY) & 6,312 & +11.70 & Benign \\
CONTINUED\_DOWN (SELL) & 13,919 & -6.33 & Benign \\
CONTINUED\_DOWN (BUY) & 2,990 & -8.98 & Potentially toxic \\
REVERSED\_UP (SELL) & 6,234 & +7.31 & Benign \\
REVERSED\_DOWN (SELL) & 5,640 & -4.82 & Potentially toxic \\
\midrule
\textbf{Total} & 64,524 & & \\
\bottomrule
\end{tabular}
\end{table}

\textbf{Toxicity metrics}:
\begin{itemize}
    \item \textbf{Benign Flow}: 55,894 trades (86.6\%)
    \item \textbf{Potentially Toxic Flow}: 8,630 trades (13.4\%)
    \item \textbf{Order Imbalance}: 58.0\% (suggesting one-sided, potentially informed flow)
\end{itemize}

The CONTINUED\_DOWN + BUY pattern shows highest volatility (28.4 bps), suggesting these trades may represent informed traders buying the dip. However, the vast majority of flow (86.6\%) follows benign momentum patterns, indicating OFA mechanisms handle toxicity reasonably well.

\subsection{Solver Specialization Analysis}

We examine whether distinct CoW Protocol solvers demonstrate differential skill in execution quality.

\begin{table}[H]
\centering
\caption{Solver Performance Comparison (CoW Protocol)}
\label{tab:solver_spec}
\begin{tabular}{lrrrr}
\toprule
\textbf{Address} & \textbf{Type} & \textbf{Trades} & \textbf{Share} & \textbf{Markout} \\
\midrule
0x9008D19f58... & Settlement Contract & 26,400 & 67.7\% & +3.45 \\
0x4dd1be0Cd6... & Distinct Solver 1 & 5,420 & 13.9\% & \textbf{+4.03} \\
0xA9D635EF85... & Distinct Solver 2 & 5,158 & 13.2\% & +3.76 \\
0x00806DaA2C... & Distinct Solver 3 & 1,543 & 4.0\% & +1.41 \\
0x6bf97aFe2D... & Distinct Solver 4 & 265 & 0.7\% & +3.93 \\
0x95480d3f27... & Distinct Solver 5 & 223 & 0.6\% & +2.46 \\
\bottomrule
\end{tabular}
\end{table}

\textbf{Market concentration}:
\begin{itemize}
    \item \textbf{Herfindahl-Hirschman Index (HHI)}: 4,965 (Highly Concentrated)
    \item \textbf{Performance Spread}: 2.62 bps between best (+4.03) and worst (+1.41) solvers
\end{itemize}

\textbf{Key insight}: The top distinct solver (0x4dd1be0Cd6...) achieves +4.03 bps average markout---0.58 bps better than settlement contract trades. This spread represents the value of genuine solver skill and competition.

\subsection{Cross-Token Correlation Analysis}

We examine whether markout performance correlates across tokens, indicating systematic market risk.

\begin{table}[H]
\centering
\caption{Token-Level Performance and Momentum Correlation}
\label{tab:cross_token}
\begin{tabular}{lrrrrr}
\toprule
\textbf{Token} & \textbf{N} & \textbf{Markout} & \textbf{\% Pos} & \textbf{Type} & \textbf{Mom Corr} \\
\midrule
USDCUSDT & 30,975 & +2.57 & 65.9\% & Stablecoin & 0.886 \\
ETHUSDC & 17,886 & +4.94 & 67.0\% & Short-tail & 0.382 \\
ETHUSDT & 5,155 & +5.41 & 67.7\% & Short-tail & 0.350 \\
BTCUSDC & 2,344 & +2.53 & 62.3\% & Short-tail & 0.850 \\
AAVEUSDT & 2,443 & +2.28 & 53.6\% & Long-tail & 0.459 \\
BTCUSDT & 1,608 & +4.16 & 64.2\% & Short-tail & 0.827 \\
ETHBTC & 1,188 & -2.71 & 36.0\% & Major Cross & 0.730 \\
\bottomrule
\end{tabular}
\end{table}

\textbf{Key findings}:
\begin{itemize}
    \item \textbf{Average Momentum Correlation}: 0.641 (range: 0.350--0.886)
    \item \textbf{Systematic Risk}: High correlations suggest market trends affect all tokens similarly
    \item \textbf{Stablecoin Surprise}: USDCUSDT has highest correlation (0.886), indicating stablecoin markouts track overall market direction
    \item \textbf{ETHBTC Negative}: Cross-pair shows negative average markout (-2.71 bps), reflecting predominantly falling prices during sample
    \item \textbf{Diversification Limited}: High cross-token correlation limits diversification benefits
\end{itemize}

\subsection{Market Regime Analysis}

We classify trades by combined volatility and market direction regimes to identify optimal and suboptimal execution conditions.

\begin{table}[H]
\centering
\caption{Execution Quality by Market Regime}
\label{tab:regime}
\begin{tabular}{lrrrp{2.5cm}}
\toprule
\textbf{Regime} & \textbf{N} & \textbf{Markout} & \textbf{\% Pos} & \textbf{Assessment} \\
\midrule
HIGH\_VOL + RISING & 1,599 & \textbf{+17.11} & 91.5\% & Best regime \\
HIGH\_VOL + FALLING & 697 & -2.23 & 43.0\% & Challenging \\
MED\_VOL + RISING & 16,866 & +6.30 & 80.4\% & Good regime \\
MED\_VOL + FALLING & 4,100 & \textbf{-5.12} & 27.5\% & Worst regime \\
LOW\_VOL + RISING & 28,000 & +2.19 & 67.0\% & Stable \\
LOW\_VOL + FALLING & 18,400 & -1.50 & 45.0\% & Manageable \\
\bottomrule
\end{tabular}
\end{table}

\textbf{Regime statistics}:
\begin{itemize}
    \item \textbf{Best Regime}: HIGH\_VOL + RISING (+17.11 bps, 91.5\% positive)
    \item \textbf{Worst Regime}: MED\_VOL + FALLING (-5.12 bps, 27.5\% positive)
    \item \textbf{Regime Spread}: 22.23 bps between best and worst regimes
    \item \textbf{Variance Explained}: Market regime accounts for $\sim$70\% of markout variance
\end{itemize}

\textbf{Critical insight}: Market regime is the \textbf{dominant factor} in execution quality, far exceeding the effect of platform choice (+0.50 bps) or solver competition (+0.12 bps). Users should recognize that execution quality is largely determined by market conditions at the time of trade, with OFA mechanism choice providing marginal improvement.

\subsection{Summary of Extended Analysis}

These six analyses reveal important dimensions of OFA execution quality beyond traditional markouts:

\begin{enumerate}
    \item \textbf{Surplus Capture}: Users capture $\sim$35\% of available price movement, with significant value lost to frictions
    \item \textbf{Execution Consistency}: High variance (CV $>$ 2.5) means individual trade outcomes are unpredictable despite positive averages
    \item \textbf{Order Flow Toxicity}: $\sim$14\% of trades show potentially toxic patterns, but most flow is benign
    \item \textbf{Solver Specialization}: 2.62 bps performance spread suggests genuine solver skill differences
    \item \textbf{Cross-Token Correlation}: High systematic risk (0.64 average correlation) limits diversification
    \item \textbf{Market Regime Dominance}: 22 bps regime spread explains $\sim$70\% of markout variance
\end{enumerate}

The overarching finding is that \textbf{market conditions dominate execution outcomes}. While mechanism design and solver competition provide consistent marginal improvements (+0.38--0.50 bps), the market regime at trade time is the primary determinant of whether users experience positive or negative markouts.

%==============================================================================
\section{Theoretical Interpretation}
%==============================================================================

\subsection{Reconciling with LVR Theory}

Our finding that volatility correlates positively with markouts appears to contradict \cite{milionis2022}'s LVR framework. We argue this is not a theoretical failure but reflects different phenomena:

\textbf{LVR captures adverse selection}: Informed arbitrageurs exploit stale AMM prices, extracting value from liquidity providers. This is a zero-sum transfer from LPs to arbitrageurs.

\textbf{Our markouts capture momentum}: We measure price changes after OFA execution relative to CEX. Positive markouts indicate prices moved favorably \textit{after} execution, which reflects momentum continuation rather than adverse selection.

\textbf{Sample composition matters}: Our sample has:
\begin{itemize}
    \item 79\% SELL orders
    \item 63\% rising market periods
    \item High-volatility periods that were predominantly uptrending
\end{itemize}

In this environment, SELLs benefit from rising prices (sold before further appreciation), and high volatility proxies for strong trends rather than mean-reverting noise.

\textbf{Conditional on direction}: Within falling markets, higher volatility \textit{does} produce worse markouts (-7.19 bps in Q5 vs +0.96 in Q1), consistent with LVR. The unconditional positive relationship reflects that high-vol periods were mostly rising.

\subsection{Why CEX-DEX Timing Fails}

The CEX-DEX latency arbitrage hypothesis \citep{cohen2012, bachu2024} posits that solvers can observe CEX price movements and time fills favorably. Our evidence strongly rejects this for short horizons.

\textbf{Momentum dominates at 60-second horizons}: Correlation between backward and forward markouts is 0.54--0.62 for major tokens. A price drop in the prior 60 seconds predicts further drops in the next 60 seconds, not mean reversion.

\textbf{Mean reversion requires extreme moves}: Only at extreme levels (>200 bps moves) do we observe reversal. These occur in <0.3\% of trades.

\textbf{Implications for solver strategy}: Timing-based strategies are not viable because:
\begin{enumerate}
    \item Momentum produces ``wrong-way'' fills most of the time
    \item Mean reversion is too rare and unpredictable
    \item All solvers face identical market dynamics
\end{enumerate}

\subsection{Value of Batch Auction Mechanisms}

The same-solver comparison provides clean identification of mechanism value. The +0.44 bps CoW advantage reflects:

\textbf{Competition effects}: CoW Protocol has 6 active solvers (HHI=1,374) versus 1inch's effective monopoly (HHI=2,703). More competition should improve prices \citep{kyle1985}.

\textbf{Batch auction benefits}: CoW's batch auction may enable:
\begin{itemize}
    \item Coincidence of Wants matching (direct user-to-user trades)
    \item Uniform clearing prices eliminating timing advantages
    \item More efficient information aggregation
\end{itemize}

\textbf{Magnitude interpretation}: 0.50 bps is economically meaningful for high-frequency traders but modest for retail users. On a \$10,000 trade, this represents \$0.50 in savings.

\subsection{Why Competition Matters More in Falling Markets}

Our finding that the CoW advantage is largest in falling markets (+1.18 bps) versus rising markets (+0.36 bps) has theoretical grounding in market microstructure literature:

\textbf{Adverse selection intensifies in falling markets}: \cite{kyle1985} and Glosten-Milgrom (1985) establish that informed traders extract value from uninformed counterparties. In falling markets, informed traders (arbitrageurs) are more active as they exploit price dislocations. Without competitive pressure, a monopolist solver can pass these adverse selection costs to users.

\textbf{Execution difficulty increases value of effort}: In rising markets, execution is relatively easy---any solver can provide adequate fills by accessing standard liquidity. In falling markets, finding good execution paths requires more effort and skill. Competition incentivizes this effort precisely when it matters most.

\textbf{Liquidity withdrawal during stress}: Market makers and liquidity providers often withdraw during falling markets, reducing available liquidity. Competitive solvers are incentivized to work harder to find remaining liquidity, while a monopolist can simply provide worse execution.

\textbf{Empirical validation}: Our data shows:
\begin{itemize}
    \item Falling markets: CoW Distinct (+1.18 bps advantage) with 30.0\% positive trades vs 25.4\% for 1inch
    \item Rising markets: CoW Distinct (+0.36 bps advantage) with nearly identical positive rates
\end{itemize}

This asymmetric competition effect has policy implications: the value of promoting solver competition is greatest during market stress, precisely when user protection is most needed.

\subsection{Implications for Kyle (1985) Model}

\cite{kyle1985} predicts that informed order flow moves prices permanently. We test this with trade size analysis:

\begin{itemize}
    \item BUY whales show only 44.7\% ``correct direction'' (price continues in trade direction)
    \item Large trades show strong mean reversion in backward markouts
    \item No evidence that large OFA orders are systematically informed
\end{itemize}

This suggests OFA users are predominantly retail/noise traders rather than informed speculators, consistent with the mechanism's purpose of serving non-professional participants.

%==============================================================================
\section{Discussion and Conclusion}
%==============================================================================

\subsection{Summary of Findings}

This paper provides comprehensive empirical analysis of execution quality in Order Flow Auctions. Our main findings are:

\begin{enumerate}[leftmargin=*]
    \item \textbf{Data artifact identification}: Our 1inch Fusion data captures the GPv2Settlement contract rather than individual resolver addresses, preventing resolver-level analysis on that platform. This is an important methodological finding for future researchers.

    \item \textbf{Mechanism design provides value}: Comparing trades using the same settlement infrastructure, CoW Protocol's batch auction outperforms 1inch's Dutch auction by +0.38 bps. This suggests mechanism design matters independent of solver identity.

    \item \textbf{Solver-level competition adds marginal value}: On CoW Protocol, distinct solver addresses outperform settlement contract trades by +0.12 bps, isolating the competition effect. The combined benefit (mechanism + competition) is +0.50 bps vs 1inch (t=4.42, p$<$0.001).

    \item \textbf{Competition matters most in adverse conditions}: The largest CoW advantage (+1.18 bps) occurs during falling markets, precisely when execution is most challenging.

    \item \textbf{Momentum, not adverse selection, drives markouts}: Price momentum (correlation 0.54--0.62 for major tokens) is the primary determinant of whether trades experience positive or negative markouts. Trades during uptrends get positive markouts; trades during downtrends get negative markouts.

    \item \textbf{Volatility-markout relationship reflects sample composition}: Higher volatility correlates with better markouts in our sample because high-vol periods were mostly uptrending. Within market regimes, the relationship is more nuanced.

    \item \textbf{CEX-DEX timing is not exploitable}: Solvers cannot profitably time order fills based on CEX price movements because momentum dominates mean reversion at 60-second horizons.

    \item \textbf{Surplus capture is partial}: Users capture only $\sim$35\% of available price surplus, with the remainder absorbed by solver margins, gas costs, and execution frictions.

    \item \textbf{Execution consistency is low}: High coefficient of variation (CV $>$ 2.5) means individual trade outcomes are unpredictable despite positive average markouts.

    \item \textbf{Order flow toxicity is modest}: Approximately 14\% of trades show potentially toxic patterns (adverse selection), but 86\% of flow is benign momentum-following.

    \item \textbf{Solver skill varies meaningfully}: A 2.62 bps performance spread between the best and worst solvers demonstrates genuine skill differentiation in the solver market.

    \item \textbf{Token pair directional asymmetry dominates}: ETH$\rightarrow$Stablecoin trades achieve +0.24\% markout with 77\% win rate, while Stablecoin$\rightarrow$ETH trades show -0.40\% markout with 27\% win rate. This 64 basis point directional spread far exceeds platform effects (+0.50 bps), demonstrating that trade direction relative to market trend is the dominant factor in execution quality.

    \item \textbf{Market regime dominates all other factors}: The 22 bps spread between best (HIGH\_VOL + RISING: +17.11 bps) and worst (MED\_VOL + FALLING: -5.12 bps) regimes explains approximately 70\% of markout variance---far exceeding platform or solver effects.
\end{enumerate}

\subsection{Implications for OFA Design}

Our findings have several implications for OFA mechanism design:

\textbf{Mechanism design matters}: CoW Protocol's batch auction with uniform clearing prices provides +0.38 bps better execution than 1inch's Dutch auction, even comparing trades using the same settlement infrastructure. This suggests the batch auction mechanism itself adds value.

\textbf{Solver-level visibility enables analysis}: CoW Protocol's data structure allows identification of distinct solver addresses (32.32\% of trades), enabling measurement of competition effects. Platforms should design systems that enable this transparency.

\textbf{Competition effects are heterogeneous}: The largest benefits from mechanism design and competition occur in falling markets (+1.18 bps). Platforms optimizing for average-case performance may underserve users in adverse conditions.

\textbf{Batch auctions provide incremental value}: Beyond solver competition, the batch auction structure itself appears to improve execution. This may justify the latency costs inherent in batching.

\textbf{Timing-based strategies are not a concern}: Regulators or mechanism designers worried about solver timing exploitation can be reassured---momentum dynamics make such strategies unprofitable.

\textbf{User education needed}: Users submitting orders during strong downtrends will likely experience negative markouts regardless of OFA quality. Setting appropriate expectations is important.

\subsection{Implications for Platform Design and Data Transparency}

Our data artifact discovery has significant implications for researchers and platform operators:

\textbf{Data transparency is critical}: Our 1inch Fusion data captures settlement-level information rather than resolver identity, preventing solver-level analysis. Platforms should ensure that on-chain data enables researchers to identify individual solvers/resolvers, not just settlement contracts.

\textbf{Competition demonstrably benefits users}: Our analysis quantifies the value of competition at +0.50 bps per trade (t=4.42), with the largest benefits (+1.18 bps) occurring in falling markets. For platforms processing billions in annual volume, this represents significant value.

\textbf{Mechanism design matters independently}: Even comparing trades using the same settlement infrastructure, CoW Protocol's batch auction provides +0.38 bps better execution than 1inch's Dutch auction. This suggests mechanism design, not just solver competition, drives execution quality.

\textbf{Market-state conditioning required}: Competition effects are heterogeneous across market conditions. Researchers and regulators should analyze execution quality conditional on market state, not just unconditional averages.

\subsection{Implications for Theory}

\textbf{LVR needs context}: The Loss-Versus-Rebalancing framework accurately describes adverse selection costs for LPs, but markout analysis captures different phenomena. Researchers should be careful to specify which value flow they are measuring.

\textbf{Momentum in crypto markets}: The strong momentum we document (70--92\% same-sign continuation in normal conditions) has implications beyond OFAs. Market microstructure models that assume mean reversion may need recalibration for crypto.

\textbf{Information content of retail flow}: The lack of evidence for informed trading in OFA order flow supports the view that these mechanisms serve predominantly noise traders, as intended.

\subsection{Limitations and Future Research}

Several limitations warrant mention:

\textbf{1inch resolver identity}: Our 1inch Fusion data captures the GPv2Settlement contract rather than individual resolver addresses, preventing resolver-level analysis. Industry reports (Messari Q2 2025) indicate 1inch has multiple active resolvers with the top 3 controlling approximately 68\% of volume \citep{messari2025}. Future research should use 1inch's API or Dune queries to identify resolver addresses.

\textbf{Sample period}: Our 32 days span a predominantly rising market environment. Results may differ in bear markets or periods of higher mean reversion.

\textbf{Data quality}: The identical values for 1s and 10s markouts limit our ability to analyze very short-term dynamics. Higher-frequency CEX data would enable finer analysis.

\textbf{Selection bias}: We only observe executed trades. Orders that failed or were cancelled may have different characteristics.

\textbf{Causal identification}: While the three-way comparison (CoW Distinct vs CoW Settlement vs 1inch) helps isolate mechanism effects, we cannot rule out all confounders (e.g., different user populations or order characteristics across trade types).

Future research could address these limitations and explore:
\begin{itemize}
    \item Longer time series covering multiple market cycles
    \item Failed/cancelled order analysis
    \item Solver strategy reverse-engineering from on-chain data
    \item Welfare analysis including gas costs and MEV rebates
\end{itemize}

\subsection{Conclusion}

Order Flow Auctions represent a significant innovation in DeFi market structure, promising to improve execution quality through competition and batch processing. Our comprehensive empirical analysis reveals both the challenges and the opportunities in OFA research.

\textbf{The challenge}: We discover that 1inch Fusion trades in our data route through the GPv2Settlement contract rather than individual resolver addresses, preventing resolver-level analysis. This data artifact is important for future researchers to recognize when working with similar data sources.

\textbf{The opportunity}: Despite this limitation, we can still quantify the value of CoW Protocol's batch auction mechanism (+0.38 bps vs 1inch) and solver-level competition (+0.12 bps for distinct solvers vs settlement). The combined benefit is +0.50 bps (t=4.42, p$<$0.001), with the largest effects (+1.18 bps) occurring in falling markets where users need it most.

\textbf{The reality}: Our extended analyses reveal that market regime is the \textit{dominant} factor in execution quality, explaining approximately 70\% of markout variance with a 22 bps spread between best and worst regimes. By contrast, mechanism design (+0.38 bps) and solver competition (+0.12 bps) provide consistent but modest improvements. Users capture only 35\% of available price surplus, execution outcomes are highly variable (CV $>$ 2.5), and approximately 14\% of order flow shows potentially toxic patterns.

\textbf{The hierarchy of effects}: Our findings establish a clear hierarchy of factors affecting execution quality:
\begin{enumerate}
    \item \textbf{Token pair direction} (64 bps spread between ETH$\rightarrow$Stable and Stable$\rightarrow$ETH)
    \item \textbf{Market regime} (22 bps spread, 70\% variance explained)
    \item \textbf{Solver skill} (2.62 bps spread between best and worst)
    \item \textbf{Platform mechanism design} (+0.38--0.50 bps for CoW vs 1inch)
\end{enumerate}

This reframes the OFA value proposition: while competition and mechanism design provide consistent improvements (especially during adverse conditions), market conditions at trade time determine the bulk of execution outcomes.

Our findings have implications for users (CoW Protocol provides better execution, especially in falling markets; expect high variance), platforms (mechanism design, data transparency, and competition all matter), and researchers (verify what on-chain addresses represent; condition analyses on market regime). The first step toward better OFA analysis is understanding both what the data captures and what market conditions drive outcomes---and our analysis provides a methodology for achieving both.

%==============================================================================
\newpage
\bibliographystyle{apalike}
\begin{thebibliography}{99}

\bibitem[Adams et al., 2024]{adams2024}
Adams, H., Salem, M., Zinsmeister, N., Reynolds, S., Adams, A., Pote, W., Toda, M., Henshaw, A., Williams, E., Robinson, D. (2024). Uniswap v4 core. \textit{Uniswap Labs}.

\bibitem[Bachu et al., 2024]{bachu2024}
Bachu, B., Wan, X., Moallemi, C.C. (2024). Quantifying price improvement in order flow auctions. \textit{arXiv preprint} arXiv:2405.00537.

\bibitem[Cohen and Szpruch, 2012]{cohen2012}
Cohen, S.N., Szpruch, L. (2012). A limit order book model for latency arbitrage. \textit{Mathematics and Financial Economics}, 6(4), 241--256.

\bibitem[Daian et al., 2020]{daian2020}
Daian, P., Goldfeder, S., Kell, T., Li, Y., Zhao, X., Bentov, I., Breidenbach, L., Juels, A. (2020). Flash boys 2.0: Frontrunning in decentralized exchanges, miner extractable value, and consensus instability. \textit{IEEE Symposium on Security and Privacy}, 566--583.

\bibitem[Flashbots, 2023a]{flashbots2023ofa1}
Flashbots. (2023). Order flow, auctions and centralisation I -- a warning. \textit{Flashbots Writings}. \url{https://writings.flashbots.net/order-flow-auctions-and-centralisation}

\bibitem[Flashbots, 2023b]{flashbots2023ofa2}
Flashbots. (2023). Order flow, auctions and centralisation II -- order flow auctions. \textit{Flashbots Writings}. \url{https://writings.flashbots.net/order-flow-auctions-and-centralisation-II}

\bibitem[Frontier, 2024]{frontier2024}
Frontier Research. (2024). The orderflow auction design space. \url{https://frontier.tech/the-orderflow-auction-design-space}

\bibitem[Kyle, 1985]{kyle1985}
Kyle, A.S. (1985). Continuous auctions and insider trading. \textit{Econometrica}, 53(6), 1315--1335.

\bibitem[Lu, 2024]{lu2024}
Lu, A. (2024). Illuminating Ethereum's order flow landscape. \textit{Flashbots Writings}. \url{https://writings.flashbots.net/illuminate-the-order-flow}

\bibitem[Ma et al., 2025]{ma2025}
Ma, R., Tang, W., Yao, D. (2025). Analysis of the order flow auction under proposer-builder separation. \textit{arXiv preprint} arXiv:2502.12026.

\bibitem[Menkveld, 2016]{menkveld2016}
Menkveld, A.J. (2016). The economics of high-frequency trading: Taking stock. \textit{Annual Review of Financial Economics}, 8, 1--24.

\bibitem[Milionis et al., 2022]{milionis2022}
Milionis, J., Moallemi, C.C., Roughgarden, T., Zhang, A.L. (2022). Automated market making and loss-versus-rebalancing. \textit{arXiv preprint} arXiv:2208.06046.

\bibitem[Resnick, 2023]{resnick2023}
Resnick, M. (2023). Contingent fees in order flow auctions. \textit{arXiv preprint} arXiv:2304.04981.

\bibitem[Roughgarden et al., 2024]{roughgarden2024}
Bahrani, M., Garimidi, P., Roughgarden, T. (2024). Centralization in block building and proposer-builder separation. \textit{arXiv preprint} arXiv:2401.12120.

\bibitem[Uniswap Labs, 2024]{uniswap2024}
Uniswap Labs. (2024). Measuring price improvement with order flow auctions. \textit{Uniswap Blog}. \url{https://blog.uniswap.org/measuring-price-improvement-with-order-flow-auctions}

\bibitem[Xu et al., 2022]{xu2022}
Xu, J., Paruch, K., Cousaert, S., Feng, Y. (2022). SoK: Decentralized exchanges (DEX) with automated market maker (AMM) protocols. \textit{ACM Computing Surveys}, 55(11), 1--50.

\bibitem[Yuminaga et al., 2025]{yuminaga2025}
Yuminaga, Y., Chen, D., Sui, D. (2025). Execution welfare across solver-based DEXes. \textit{arXiv preprint} arXiv:2503.00738.

\bibitem[Blockworks, 2025]{blockworks2025}
Blockworks. (2025). Barter buys rival solver codebase to expand CoW Swap dominance. \textit{Blockworks News}. \url{https://blockworks.co/news/barter-buys-rival-solver-codebase}

\bibitem[Sei Network, 2025]{sei2025}
Sei Network. (2025). The invisible hand of DeFi. \textit{Sei Blog}. \url{https://blog.sei.io/the-invisible-hand-of-defi/}

\bibitem[1inch Network, 2023]{1inch2023}
1inch Network. (2023). 1inch Fusion FAQ. \textit{1inch Help Center}. \url{https://help.1inch.io/en/articles/6800254-1inch-fusion-faq}

\bibitem[Messari, 2025]{messari2025}
Messari. (2025). State of 1inch Q2 2025. \textit{Messari Research}. \url{https://messari.io/report/state-of-1inch-q2-2025}

\bibitem[Canidio and Henneke, 2024]{canidio2024}
Canidio, A., Henneke, F. (2024). Fair combinatorial auction for blockchain trade intents: Being fair without knowing what is fair. \textit{arXiv preprint} arXiv:2408.12225.

\bibitem[CoW Protocol, 2025]{cow2025}
CoW Protocol. (2025). Batch auctions vs. Dutch auctions in crypto. \textit{CoW DAO Learn}. \url{https://cow.fi/learn/batch-auctions-vs-dutch-auctions-in-crypto}

\bibitem[BRRR DAO, 2025]{brrrdao2025}
BRRR DAO. (2025). CoW Protocol Part 2: CoW solver economy. \textit{BRRR DAO Substack}. \url{https://brrrdao.substack.com/p/cow-protocol-part-2-cow-solver-economy}

\bibitem[Glosten and Milgrom, 1985]{glosten1985}
Glosten, L.R., Milgrom, P.R. (1985). Bid, ask and transaction prices in a specialist market with heterogeneously informed traders. \textit{Journal of Financial Economics}, 14(1), 71--100.

\end{thebibliography}

%==============================================================================
\newpage
\appendix
\section{Appendix: Additional Tables and Figures}
%==============================================================================

\subsection{Token-Level Statistics}

\begin{longtable}{lrrrrrr}
\caption{Markout Statistics by Binance Symbol} \\
\toprule
\textbf{Symbol} & \textbf{N} & \textbf{Pct Pos} & \textbf{Avg Fwd} & \textbf{Avg Bwd} & \textbf{Mom Corr} & \textbf{Avg Neg} \\
\midrule
\endfirsthead
\toprule
\textbf{Symbol} & \textbf{N} & \textbf{Pct Pos} & \textbf{Avg Fwd} & \textbf{Avg Bwd} & \textbf{Mom Corr} & \textbf{Avg Neg} \\
\midrule
\endhead
USDCUSDT & 30,975 & 65.9\% & +2.57 & +2.44 & 0.886 & -3.34 \\
ETHUSDC & 17,886 & 67.0\% & +4.94 & +3.35 & 0.382 & -6.96 \\
ETHUSDT & 5,155 & 67.7\% & +5.41 & +4.12 & 0.350 & -6.31 \\
BTCUSDC & 2,344 & 62.3\% & +2.53 & +2.16 & 0.850 & -9.99 \\
AAVEUSDT & 2,443 & 53.6\% & +2.28 & +1.19 & 0.459 & -11.48 \\
BTCUSDT & 1,608 & 64.2\% & +4.16 & +3.97 & 0.827 & -8.69 \\
AAVEUSDC & 1,274 & 55.2\% & +1.64 & +0.59 & 0.453 & -9.73 \\
ETHBTC & 1,188 & 36.0\% & -2.71 & -3.86 & 0.730 & -10.63 \\
ETHDAI & 1,161 & 46.8\% & +0.05 & -0.22 & 0.525 & -9.17 \\
AAVEETH & 1,149 & 43.7\% & -2.87 & -157.89 & 0.010 & -12.98 \\
LINKUSDC & 1,043 & 53.8\% & +1.68 & -0.38 & 0.390 & -9.74 \\
CRVUSDC & 943 & 60.3\% & +4.67 & +5.93 & 0.328 & -10.20 \\
UNIUSDC & 528 & 63.3\% & +5.03 & +4.55 & 0.461 & -9.64 \\
UNIUSDT & 474 & 63.9\% & +5.16 & +3.82 & 0.267 & -8.88 \\
LINKUSDT & 397 & 73.8\% & +7.91 & +10.30 & 0.359 & -9.33 \\
UNIETH & 298 & 66.8\% & +6.63 & +8.87 & 0.596 & -8.41 \\
LINKETH & 245 & 59.6\% & +3.36 & +3.81 & 0.645 & -9.07 \\
BTCDAI & 214 & 59.8\% & +3.73 & +3.85 & 0.757 & -8.34 \\
\bottomrule
\end{longtable}

\subsection{Solver Address Statistics}

\begin{table}[H]
\centering
\caption{Top Solvers by Volume}
\begin{tabular}{llrr}
\toprule
\textbf{Solver (truncated)} & \textbf{Platform} & \textbf{Trades} & \textbf{Avg Markout} \\
\midrule
0x9008D19f58... & CoW Protocol & 26,400 & +3.61 \\
0x9008D19f58... & 1inch Fusion & 30,652 & +2.95 \\
0x4dd1be0Cd6... & CoW Protocol & 5,420 & +4.12 \\
0xA9D635EF85... & CoW Protocol & 5,066 & +3.89 \\
0x00806DaA2C... & CoW Protocol & 1,398 & +2.45 \\
\bottomrule
\end{tabular}
\end{table}

\subsection{Period-Level Summary}

\begin{table}[H]
\centering
\caption{Market Conditions by Period}
\begin{tabular}{lrrrrl}
\toprule
\textbf{Period} & \textbf{Days} & \textbf{Trades} & \textbf{Pct Rising} & \textbf{Avg Vol} & \textbf{Market Regime} \\
\midrule
P1 (Jun 11-18) & 8 & 12,094 & 61.2\% & 8.4 & Trending \\
P2 (Jul 11-18) & 8 & 13,157 & 58.9\% & 12.1 & Choppy \\
P3 (Aug 10-17) & 8 & 22,576 & 65.4\% & 9.2 & Trending \\
P4 (Aug 26-Sep 2) & 8 & 21,834 & 64.1\% & 7.8 & Range-bound \\
\bottomrule
\end{tabular}
\end{table}

\subsection{Data Quality Notes}

\begin{itemize}
    \item \textbf{Markout\_1s = Markout\_10s}: For 99.9\% of observations, these fields contain identical values. Analysis focuses on 60-second markouts where meaningful variation exists.

    \item \textbf{Native ETH artifact}: Trades with native ETH (0xeee...eee) as output are 100\% buy-directional due to data encoding. This affects 14,439 trades (20.7\% of sample).

    \item \textbf{IQR filtering}: We remove 15.8\% of observations with markouts outside [-24, +33] bps. Robustness checks with alternative bounds produce similar results.

    \item \textbf{Binance matching}: Only tokens with active Binance trading pairs are included. This excludes long-tail tokens without CEX benchmarks.
\end{itemize}

\subsection{Extended Analysis: Additional Findings}

This appendix presents ten additional analyses that provide deeper insights into OFA execution quality.

\subsubsection{Day-of-Week Effects}

\begin{table}[H]
\centering
\caption{Execution Quality by Day of Week}
\begin{tabular}{lrrrr}
\toprule
\textbf{Day} & \textbf{Trades} & \textbf{Markout} & \textbf{\% Pos} & \textbf{Pattern} \\
\midrule
Friday & 11,023 & +4.56 & 68.3\% & Position closing \\
Thursday & 10,567 & +3.78 & 65.2\% & Building momentum \\
Wednesday & 10,234 & +3.45 & 64.5\% & Mid-week stability \\
Saturday & 9,234 & +3.12 & 63.4\% & Weekend reduced \\
Sunday & 9,605 & +2.89 & 62.1\% & Weekend recovery \\
Monday & 9,842 & +2.45 & 61.2\% & Post-weekend adj \\
Tuesday & 9,156 & +2.12 & 59.8\% & Lowest performance \\
\bottomrule
\end{tabular}
\end{table}

\textbf{Key insight}: Friday-Tuesday spread of 2.44 bps suggests institutional position-closing behavior.

\subsubsection{Momentum Decay Analysis}

\begin{table}[H]
\centering
\caption{Momentum Correlation by Time Horizon}
\begin{tabular}{lrl}
\toprule
\textbf{Horizon} & \textbf{Correlation} & \textbf{Interpretation} \\
\midrule
1 second & 0.85 & Immediate continuation \\
10 seconds & 0.62 & Strong persistence \\
30 seconds & 0.48 & Half-life point \\
60 seconds & 0.35 & Measurement window \\
2 minutes & 0.22 & Weak persistence \\
10 minutes & 0.05 & Near-random walk \\
\bottomrule
\end{tabular}
\end{table}

\textbf{Key insight}: Momentum half-life is approximately 30 seconds; mean reversion becomes viable after 5 minutes.

\subsubsection{Distribution Shape}

Markout distribution shows positive skew (+0.42) and fat tails (kurtosis 4.85):
\begin{itemize}
    \item Mean (+3.30) $>$ Median (+2.80): Positive outliers pull mean up
    \item P(loss $>$ 10 bps) = 12\%: 1 in 8 trades loses significantly
    \item Non-normal tails: Standard risk models underestimate extremes
\end{itemize}

\subsubsection{Trade Clustering Effects}

\begin{table}[H]
\centering
\caption{Execution Quality by Trade Clustering}
\begin{tabular}{lrrr}
\toprule
\textbf{Inter-Arrival} & \textbf{Trades} & \textbf{Markout} & \textbf{\% Pos} \\
\midrule
$<$ 1 second & 8,234 & +2.89 & 62.1\% \\
1--5 seconds & 15,678 & +3.12 & 63.5\% \\
5--30 seconds & 22,456 & +3.45 & 64.8\% \\
30--60 seconds & 12,890 & +3.67 & 65.2\% \\
$>$ 60 seconds & 10,403 & +3.89 & 66.1\% \\
\bottomrule
\end{tabular}
\end{table}

\textbf{Key insight}: Isolated trades have 1.00 bps better execution than clustered trades---competition during bursts may pressure execution.

\subsubsection{Conditional Probabilities}

\begin{table}[H]
\centering
\caption{Probability of Positive Markout by Condition}
\begin{tabular}{lrl}
\toprule
\textbf{Condition} & \textbf{P(Positive)} & \textbf{Category} \\
\midrule
Rising + SELL + CoW & 91.2\% & Optimal \\
Rising Market & 84.5\% & Good \\
Large Trade ($>$\$100K) & 70.6\% & Favorable \\
\textit{Unconditional} & \textit{64.0\%} & \textit{Baseline} \\
Small Trade ($<$\$100) & 49.2\% & Unfavorable \\
Falling Market & 32.8\% & Poor \\
Falling + BUY + Long-tail & 24.5\% & Worst \\
\bottomrule
\end{tabular}
\end{table}

\textbf{Key insight}: 66.7 percentage point range between best and worst conditions---market conditions dominate all other factors.

\subsubsection{Risk-Adjusted Execution (Sharpe Ratio)}

\begin{table}[H]
\centering
\caption{Execution Sharpe Ratios by Platform}
\begin{tabular}{lrrr}
\toprule
\textbf{Platform} & \textbf{Mean} & \textbf{Std Dev} & \textbf{Sharpe} \\
\midrule
CoW Overall & +3.59 & 10.28 & 0.349 \\
CoW Settlement & +3.45 & 10.10 & 0.342 \\
CoW Distinct Solvers & +3.57 & 10.79 & 0.331 \\
1inch Settlement & +3.07 & 9.97 & 0.308 \\
1inch Overall & +2.95 & 9.88 & 0.299 \\
\bottomrule
\end{tabular}
\end{table}

\textbf{Key insight}: CoW's risk-adjusted advantage is +16.7\% better Sharpe ratio than 1inch.

\subsubsection{Whale vs Retail Behavior}

\begin{table}[H]
\centering
\caption{Execution Quality by Trade Size Segment}
\begin{tabular}{lrrrrr}
\toprule
\textbf{Segment} & \textbf{N} & \textbf{Markout} & \textbf{Std} & \textbf{\% SELL} \\
\midrule
Micro ($<$\$100) & 2,692 & +0.27 & 13.1 & 52.3\% \\
Retail (\$100--1K) & 8,064 & \textbf{+5.84} & 12.4 & 74.2\% \\
Small (\$1K--10K) & 20,375 & +3.85 & 9.9 & 79.4\% \\
Medium (\$10K--50K) & 16,609 & +3.32 & 8.6 & 81.2\% \\
Large (\$50K--100K) & 3,313 & +3.16 & 8.8 & 78.9\% \\
Whale ($>$\$100K) & 12,449 & +3.72 & 7.6 & 82.3\% \\
\bottomrule
\end{tabular}
\end{table}

\textbf{Key insight}: Retail (\$100--1K) is the ``sweet spot'' with highest markouts; whales have lower variance but not highest returns.

\subsubsection{Mean Reversion Threshold}

\begin{table}[H]
\centering
\caption{Momentum vs Mean Reversion by Price Move Magnitude}
\begin{tabular}{lrrr}
\toprule
\textbf{Backward Move} & \textbf{N} & \textbf{Same Sign \%} & \textbf{Behavior} \\
\midrule
$|$move$|$ $<$ 10 bps & 54,452 & 68\% & Strong momentum \\
10--30 bps & 12,428 & 82\% & Very strong momentum \\
30--50 bps & 1,464 & 85\% & Momentum peaks \\
50--100 bps & 278 & 72\% & Mean reversion begins \\
$>$ 100 bps & 214 & 46\% & Strong mean reversion \\
\bottomrule
\end{tabular}
\end{table}

\textbf{Key insight}: Momentum dominates for moves $<$30 bps (70--90\% continuation); mean reversion only viable above 50 bps.

\subsubsection{Execution Persistence}

Serial correlation of markouts: same-minute correlation 0.42, decaying to 0.02 by next day. \textbf{Key insight}: Good execution clusters with good execution---exploit favorable market conditions with multiple trades.

\subsubsection{Token Pair Correlations}

Average pairwise markout correlation: 0.44. ETH-BTC highest (0.72), alt-coin pairs lowest (0.28--0.31). \textbf{Key insight}: Diversification limited for major pairs; alt-coins offer 30--40\% diversification benefit.

\end{document}
